%% kapa: exam=1 round=37 year=2026 period=1 form=A subject=civil_law subject_ko=민법 seq=29 number=29
%% source: Q-Net 감정평가사 기출문제 (cst00309) artlSeq=5258924 file=raw/1008/5258924/2026년 제37회 감정평가사 1차 1교시.pdf
%% conversion: kordoc 4.24.51 -> tools/gen1.py (auto); review=proofread-sonnet; flags=none
\begin{question}{29}
등기부취득시효에 관한 설명으로 옳지 않은 것은? (다툼이 있으면 판례에 따름)
\begin{choices}
  \item 등기부취득시효의 요건으로서의 소유자로 등기한 자라 함은 적법ㆍ유효한 등기를 마친 자일 필요는 없고 중복등기가 아닌 한 무효인 등기라도 무방하다.
  \item 등기부시효취득에 있어서 무과실은 점유취득에 관한 것이고 등기에 관한 것이 아니다.
  \item 시효완성자의 선의ㆍ무과실은 시효기간 내내 계속되어야 하는 것은 아니다.
  \item 부동산의 소유자로 등기된 기간과 점유기간이 때를 같이 하여 10년을 경과하여야만 시효취득이 인정된다.
  \item 부동산 매수인이 매도인의 부동산 처분권한을 조사했더라면 그 처분권한이 없음을 알 수 있었음에도 이를 조사하지 않은 경우, 매수인의 등기부취득시효는 완성되지 않는다.
\end{choices}
\end{question}
\answer{4}
